% !PRINTED 360
\weberpage{360}
%<<CONT
fremden Zahlen bedeutet und $q$ in dem Produkt $\Pi$ die voneinander
 verschieden in $Q$ aufgehenden Primzahlen durchläuft.

Bezeichnen wir mit
\[
\varrho_1, \varrho_2, \ldots, \varrho_\nu
\]
ein volles Repräsentantensystem der zu $Q$ teilerfremden Zahlen
in $\Omega$, nach dem Modul $Q$ genommen, so ist in gleicher Weise
\weq{6}{O = \varrho_1 O_0 + \varrho_2 O_0 + \cdots + \varrho_\nu O_0}
und
\weq{7}{\nu = \psi(Q) = (O, O_0)}
ist der Index des Teilers $O_0$ in bezug auf $O$.

Die Zahl $\psi(Q)$ haben wir in § 168, (3) des II. Bandes allgemein
 bestimmt, auch für den Fall, daß an Stelle von $Q$ ein
Ideal tritt. Da nun hier $N(Q) = Q^2$ ist, so haben wir:
\weq{8}{\psi(Q) = Q^2\, \Pi\left(1 - \frac{1}{N(\mathfrak{q})}\right),}
wenn $\mathfrak{q}$ die voneinander verschiedenen idealen Primteiler von $Q$
durchläuft.

Nun ist (§ 92)
\[
\begin{aligned}
&\text{1. } N(\mathfrak{q}) = q, \quad \text{wenn } (\Delta,q) = 0,\\
&\text{2. } N(\mathfrak{q}) = q, \quad \text{wenn } (\Delta,q) = +1,\\
&\text{3. } N(\mathfrak{q}) = q^2, \quad \text{wenn } (\Delta,q) = -1.
\end{aligned}
\]
Wir wollen diese dreierlei Primzahlen mit $q_1$, $q_2$, $q_3$ bezeichnen
 und bemerken, daß jede Primzahl $q_2$ die Norm von
zwei verschiedenen Primidealen $\mathfrak{q}$ ist. Demnach ergibt sich aus (8)
\weq{9}{\psi(Q) = Q^2\, \Pi\left(1 - \frac{1}{q_1}\right) \Pi\left(1 - \frac{1}{q_2}\right)^2 \Pi\left(1 - \frac{1}{q_3^2}\right).}

Nun ist der Index von $O'$ in bezug auf $O$ nach Bd.\,II, § 2 (4):
\[
j = \frac{\nu}{\mu} = (O, O') = \frac{(O, O_0)}{(O', O_0)} = \frac{\psi(Q)}{\varphi(Q)}
\]
und
\[
\varphi(Q) = Q\, \Pi\left(1 - \frac{1}{q_1}\right) \Pi\left(1 - \frac{1}{q_2}\right) \Pi\left(1 - \frac{1}{q_3}\right);
\]
folglich ergibt sich:
\[
(O, O') = Q\, \Pi\left(1 - \frac{1}{q_2}\right) \Pi\left(1 + \frac{1}{q_3}\right),
\]
wofür dann auch gesetzt werden kann (nach 1., 2., 3.):
\weq{10}{j = (O, O') = Q\, \Pi\left(1 - \frac{(\Delta,q)}{q}\right).}
